Las figuras de este anejo no se emplean como prueba cuantitativa principal. Su
función es mostrar el tipo de evidencia visual que se inspecciona después de una
simulación cerrada y dar contexto a las métricas agregadas del capítulo de
resultados. Se han escogido casos limpios, casos con perturbaciones y
trayectorias con cambios geométricos para que el lector pueda ver el alcance
práctico del banco de simulación.

\subsubsection*{Atlas de trayectorias}

Antes de entrar en señales temporales, la \autoref{fig:atlas-trayectorias-id},
la \autoref{fig:atlas-trayectorias-ood} y la
\autoref{fig:atlas-trayectoria-helix-3d} sitúan visualmente la variedad de
referencias. Las dos primeras figuras resumen familias conocidas y trayectorias
fuera de distribución, mientras que la tercera separa una hélice para que se
aprecie su componente tridimensional.

\begin{figure}[H]
\centering
\includegraphics[width=\textwidth]{atlas_trayectorias_id.pdf}
\caption{Muestra visual de familias conocidas: \texttt{hold}, \texttt{circle},
\texttt{lissajous} y \texttt{waypoint}. Fuente: elaboración propia.}
\label{fig:atlas-trayectorias-id}
\end{figure}

\begin{figure}[H]
\centering
\includegraphics[width=\textwidth]{atlas_trayectorias_ood.pdf}
\caption{Muestra visual de trayectorias fuera de distribución:
\texttt{lemniscate}, \texttt{lissajous} 3D, composición y hélice. Fuente:
elaboración propia.}
\label{fig:atlas-trayectorias-ood}
\end{figure}

\begin{figure}[H]
\centering
\includegraphics[width=0.78\textwidth]{atlas_trayectoria_helix_3d.pdf}
\caption{Vista tridimensional de una hélice OOD representativa. Fuente:
elaboración propia.}
\label{fig:atlas-trayectoria-helix-3d}
\end{figure}

\subsubsection*{Episodios seleccionados}

La \autoref{tab:casos-telemetria-anejo} resume los episodios incluidos en las
figuras siguientes. La selección no busca maximizar o minimizar una métrica,
sino cubrir situaciones distintas: estacionario, seguimiento curvo limpio,
seguimiento con ruido y viento, navegación por waypoints y una trayectoria
compuesta que encadena segmentos.

\begin{table}[H]
\caption{Casos de telemetría incluidos en el anejo.}
\label{tab:casos-telemetria-anejo}
\small
\begin{tabularx}{\textwidth}{p{0.24\textwidth} p{0.22\textwidth} X}
\toprule
Caso & Escenario & Motivo de inclusión \\
\midrule
\texttt{hover\_clean} & Estacionario limpio & Verifica equilibrio vertical,
empuje próximo al peso y ausencia de excitaciones laterales. \\
\texttt{lissajous\_clean} & Curva 3D limpia & Muestra seguimiento continuo con
acoplamiento simultáneo en los tres ejes. \\
\texttt{circle\_noisy\_wind} & Círculo con drag, ruido y viento & Ilustra la
respuesta ante perturbaciones y límites de actuador más restrictivos. \\
\texttt{waypoint\_clean} & Secuencia por puntos & Enseña cambios discretos de
referencia, tramos de avance y parada en destino. \\
\texttt{composite\_ood} & Trayectoria compuesta & Presenta transiciones entre
familias y mayor diversidad de consignas dentro de un mismo episodio. \\
\bottomrule
\end{tabularx}
\end{table}

\subsubsection*{Lectura geométrica}

La \autoref{fig:sim-telemetry-geometry-a} y la
\autoref{fig:sim-telemetry-geometry-b} comparan la referencia y el estado en el
plano horizontal para cuatro situaciones representativas. En hover la
geometría queda prácticamente degenerada en un punto; en Lissajous y círculo
aparece el seguimiento curvo; y en waypoint se observan los tramos rectos con
cambios de dirección.

\begin{figure}[H]
\centering
\begin{subfigure}[t]{0.48\textwidth}
\centering
\includegraphics[width=\textwidth]{sim_hover_clean/trajectory_xy.pdf}
\caption{\texttt{hover\_clean}.}
\end{subfigure}\hfill
\begin{subfigure}[t]{0.48\textwidth}
\centering
\includegraphics[width=\textwidth]{sim_lissajous_clean/trajectory_xy.pdf}
\caption{\texttt{lissajous\_clean}.}
\end{subfigure}
\caption{Seguimiento geométrico en escenarios limpios. Fuente: elaboración
propia.}
\label{fig:sim-telemetry-geometry-a}
\end{figure}

\begin{figure}[H]
\centering
\begin{subfigure}[t]{0.48\textwidth}
\centering
\includegraphics[width=\textwidth]{sim_circle_noisy_wind/trajectory_xy.pdf}
\caption{\texttt{circle\_noisy\_wind}.}
\end{subfigure}\hfill
\begin{subfigure}[t]{0.48\textwidth}
\centering
\includegraphics[width=\textwidth]{sim_waypoint_clean/trajectory_xy.pdf}
\caption{\texttt{waypoint\_clean}.}
\end{subfigure}
\caption{Seguimiento geométrico en un caso perturbado y en una misión por
waypoints. Fuente: elaboración propia.}
\label{fig:sim-telemetry-geometry-b}
\end{figure}

\subsubsection*{Error y actitud}

Para interpretar una simulación no basta con ver la trayectoria. La
\autoref{fig:sim-telemetry-error-attitude} combina error de seguimiento y
actitud en dos episodios: un círculo con viento, ruido y drag, y una trayectoria
compuesta OOD. Esto permite relacionar errores transitorios con los cambios de
actitud que el lazo interno debe estabilizar.

\begin{figure}[H]
\centering
\begin{subfigure}[t]{0.48\textwidth}
\centering
\includegraphics[width=\textwidth]{sim_circle_noisy_wind/tracking_error.pdf}
\caption{Error en \texttt{circle\_noisy\_wind}.}
\end{subfigure}\hfill
\begin{subfigure}[t]{0.48\textwidth}
\centering
\includegraphics[width=\textwidth]{sim_circle_noisy_wind/attitude_time.pdf}
\caption{Actitud en \texttt{circle\_noisy\_wind}.}
\end{subfigure}

\vspace{0.5em}
\begin{subfigure}[t]{0.48\textwidth}
\centering
\includegraphics[width=\textwidth]{sim_composite_ood/tracking_error.pdf}
\caption{Error en \texttt{composite\_ood}.}
\end{subfigure}\hfill
\begin{subfigure}[t]{0.48\textwidth}
\centering
\includegraphics[width=\textwidth]{sim_composite_ood/attitude_time.pdf}
\caption{Actitud en \texttt{composite\_ood}.}
\end{subfigure}
\caption{Relación visual entre error de seguimiento y actitud en dos episodios
con mayor exigencia dinámica. Fuente: elaboración propia.}
\label{fig:sim-telemetry-error-attitude}
\end{figure}

\subsubsection*{Series temporales por eje}

La lectura por ejes de la \autoref{fig:sim-telemetry-position-time} resulta
útil cuando la geometría en planta no muestra toda la exigencia de la misión.
En Lissajous se aprecia la variación simultánea de altura, mientras que en
waypoint se distinguen tramos de avance y mesetas asociadas a la llegada a cada
punto.

\begin{figure}[H]
\centering
\begin{subfigure}[t]{0.48\textwidth}
\centering
\includegraphics[width=\textwidth]{sim_lissajous_clean/position_time.pdf}
\caption{Posición por eje en \texttt{lissajous\_clean}.}
\end{subfigure}\hfill
\begin{subfigure}[t]{0.48\textwidth}
\centering
\includegraphics[width=\textwidth]{sim_waypoint_clean/position_time.pdf}
\caption{Posición por eje en \texttt{waypoint\_clean}.}
\end{subfigure}
\caption{Evolución temporal de la posición frente a la referencia en dos
familias con lectura geométrica diferente. Fuente: elaboración propia.}
\label{fig:sim-telemetry-position-time}
\end{figure}

\subsubsection*{Actuación}

La \autoref{fig:sim-telemetry-actuation} muestra señales de actuación en tres
casos. El hover permite reconocer el nivel de empuje cercano al peso; el caso
con viento exige variaciones sostenidas de control; y la trayectoria compuesta
presenta cambios de consigna asociados a sus segmentos. Estas figuras ayudan a
ver que la comparación no se limita a posiciones, sino que también se vigilan
empujes, momentos y velocidades de rotor.

\begin{figure}[H]
\centering
\begin{subfigure}[t]{0.48\textwidth}
\centering
\includegraphics[width=\textwidth]{sim_hover_clean/control_effort.pdf}
\caption{Esfuerzo de control en \texttt{hover\_clean}.}
\end{subfigure}\hfill
\begin{subfigure}[t]{0.48\textwidth}
\centering
\includegraphics[width=\textwidth]{sim_circle_noisy_wind/control_effort.pdf}
\caption{Esfuerzo de control en \texttt{circle\_noisy\_wind}.}
\end{subfigure}

\vspace{0.5em}
\begin{subfigure}[t]{0.48\textwidth}
\centering
\includegraphics[width=\textwidth]{sim_composite_ood/control_effort.pdf}
\caption{Esfuerzo de control en \texttt{composite\_ood}.}
\end{subfigure}\hfill
\begin{subfigure}[t]{0.48\textwidth}
\centering
\includegraphics[width=\textwidth]{sim_composite_ood/rotor_speeds.pdf}
\caption{Velocidad de rotores en \texttt{composite\_ood}.}
\end{subfigure}
\caption{Señales de actuación inspeccionadas tras la simulación cerrada.
Fuente: elaboración propia.}
\label{fig:sim-telemetry-actuation}
\end{figure}

Finalmente, la \autoref{fig:sim-telemetry-rates} completa la lectura de
actuación con velocidades angulares del cuerpo. Estas señales permiten detectar
si una misión aparentemente aceptable en posición se apoya en oscilaciones de
actitud o en demandas rápidas del lazo interno.

\begin{figure}[H]
\centering
\begin{subfigure}[t]{0.48\textwidth}
\centering
\includegraphics[width=\textwidth]{sim_circle_noisy_wind/angular_velocity_time.pdf}
\caption{Velocidades angulares en \texttt{circle\_noisy\_wind}.}
\end{subfigure}\hfill
\begin{subfigure}[t]{0.48\textwidth}
\centering
\includegraphics[width=\textwidth]{sim_waypoint_clean/angular_velocity_time.pdf}
\caption{Velocidades angulares en \texttt{waypoint\_clean}.}
\end{subfigure}
\caption{Velocidades angulares del vehículo en cuerpo para dos escenarios de
seguimiento. Fuente: elaboración propia.}
\label{fig:sim-telemetry-rates}
\end{figure}
